% =====================================================================
% 08-Appendix-Misc.tex
% Geometric Hallucination Taxonomy — Mathematical Spine
% Borealis S. Hedling — 2026
% =====================================================================

\section*{08. Appendix: Miscellaneous Lemmas and Numerical Notes}

This appendix contains auxiliary results, discretization formulas,
approximation lemmas, and supporting material used throughout the
mathematical spine. These results are not operator-specific but are
required for numerical validation, Python toolkit alignment, and Lean
formalization consistency.

% ---------------------------------------------------------------------
\subsection*{8.1 Jacobi Field Approximation}

Let $\gamma$ be a geodesic and $J$ a Jacobi field along $\gamma$:
\[
\nabla_{\dot{\gamma}}^2 J + R(J,\dot{\gamma})\dot{\gamma} = 0.
\]

In discrete settings, the second covariant derivative is approximated by
\[
\nabla_{\dot{\gamma}}^2 J(t)
\approx
\frac{J(t+1) - 2J(t) + J(t-1)}{\Delta t^2}.
\]

This approximation is used in curvature-based diagnostics.

% ---------------------------------------------------------------------
\subsection*{8.2 Discrete Christoffel Symbols}

Given a coordinate chart, Christoffel symbols are approximated by
finite differences:
\[
\Gamma^i_{\;jk}(x_t)
\approx
\frac{1}{2} g^{im}(x_t)
\left(
\frac{g_{km}(x_{t+1}) - g_{km}(x_{t-1})}{2\Delta t}
+
\frac{g_{jm}(x_{t+1}) - g_{jm}(x_{t-1})}{2\Delta t}
-
\frac{g_{jk}(x_{t+1}) - g_{jk}(x_{t-1})}{2\Delta t}
\right).
\]

This is the discrete version used in curvature and holonomy modules.

% ---------------------------------------------------------------------
\subsection*{8.3 Parallel Transport Approximation}

Parallel transport of a vector $V_t$ along a discrete trajectory is
approximated by
\[
V_{t+1}
= V_t - \Gamma(x_t)(\dot{x}_t, V_t)\Delta t.
\]

This formula is used in holonomy drift experiments.

% ---------------------------------------------------------------------
\subsection*{8.4 Discrete Eigenvalue Drift}

Given matrices $A_t$ and $A_{t+1}$, eigenvalue drift is computed by
\[
\Delta_\lambda(t)
= \lambda(t+1) - \lambda(t),
\]
where eigenvalues are obtained via numerical eigendecomposition.

This approximation is used in spectral cascade diagnostics.

% ---------------------------------------------------------------------
\subsection*{8.5 Metric Determinant Approximation}

The determinant of the metric matrix is computed numerically via
\[
\det(g_{ij}(x_t))
= \prod_{i=1}^n \sigma_i,
\]
where $\sigma_i$ are singular values of $g_{ij}(x_t)$.

This is used in topological collapse detection.

% ---------------------------------------------------------------------
\subsection*{8.6 Fusion–Field Numerical Stability}

Given latent vectors $x_t$ and $y_t$, the fusion–field output is
\[
F_t = \alpha x_t + \beta y_t + \gamma \sigma(x_t, y_t).
\]

To avoid numerical overflow, nonlinear terms are clipped:
\[
\sigma(x_t, y_t)
\leftarrow
\min\left( \sigma(x_t, y_t), \sigma_{\max} \right).
\]

This ensures stable numerical evaluation of fusion spikes.

% ---------------------------------------------------------------------
\subsection*{8.7 Auxiliary Lemma: Threshold Crossing}

\begin{lemma}[Threshold Crossing]
Let $z(t)$ be a diagnostic quantity (curvature, spectral derivative,
holonomy norm, metric rank, or fusion norm). If $z(t)$ is continuous in
$t$ and $z(t_0) > z_{\mathrm{thr}}$, then there exists $\delta > 0$ such
that $z(t) > z_{\mathrm{thr}}$ for all $t \in (t_0 - \delta, t_0 + \delta)$.
\end{lemma}

\begin{proof}[Sketch]
Continuity ensures that values near $t_0$ remain above the threshold.
\end{proof}

% ---------------------------------------------------------------------
\subsection*{8.8 Auxiliary Lemma: Instability Propagation}

\begin{lemma}[Propagation]
If operator $A$ is unstable at time $t_0$ and operator $B$ depends
smoothly on $A$, then $B$ is unstable at some time $t_1$ near $t_0$.
\end{lemma}

\begin{proof}[Sketch]
Smooth dependence ensures that instability in $A$ induces instability in
$B$ through the interaction map $\mathcal{I}$.
\end{proof}

% ---------------------------------------------------------------------
\subsection*{8.9 Numerical Stability Conditions}

To ensure numerical stability in experiments:
\begin{itemize}
\item step size $\Delta t$ must satisfy $\Delta t < \Delta t_{\max}$,
\item metric matrices must be regularized to avoid near-degeneracy,
\item eigenvalue computations must use stable algorithms,
\item nonlinear fusion terms must be clipped or normalized.
\end{itemize}

These conditions ensure reproducibility of diagnostic results.

% ---------------------------------------------------------------------
\subsection*{8.10 Lean Formalization Notes}

For mechanization in Lean:
\begin{itemize}
\item discrete approximations correspond to sequences,
\item continuous operators correspond to smooth structures,
\item thresholds correspond to fixed constants,
\item instability predicates correspond to decidable propositions.
\end{itemize}

These notes ensure consistency between the `.tex` spine and Lean files.

% =====================================================================
% End of 08-Appendix-Misc.tex
% =====================================================================
